\documentclass{exam}

\newif\ifA
\newif\ifkey
\newif\ifprintselected   % required by exam_config's \ansbox

%%%%%%%%%%%%%%%%%%%%%%%%%%%%%%%%%%%%%%%%%%%%%%%%%%%%%%
%%% SET THESE IF VARIABLES
%%%%%%%%%%%%%%%%%%%%%%%%%%%%%%%%%%%%%%%%%%%%%%%%%%%%%%
% Lab name -- labs are named by the lesson they belong to, never by a
% sequential lab number (Neil, 2026-10-12): this is the L21 Lab. The packet replaces the old L21 Lab sheet and folds in
% the retired L21Lab practice set (Pre-Lab and Writeup), 2026-10-21.
\def\labLesson{21}

% is this the key or not?
\ifdefined\iskey
	\ifnum\iskey=1
		\keytrue
	\else
		\keyfalse
	\fi
\else
	%% Pick one manually
	\keyfalse % if this is NOT the key
%	\keytrue     % if this IS the key
\fi

% set up the header and footer - change for every term
\ifdefined\thisTermIn
	\def\thisterm{\thisTermIn}
\else
	\def\thisterm{Fall 2026}
\fi
%
\def\thisclass{ECE 444}

%%%%%%%%%%%%%%%%%%%%%%%%%%%%%%%%%%%%%%%%%%%%%%%%%%%%%%
%%% END SET THESE IF VARIABLES
%%%%%%%%%%%%%%%%%%%%%%%%%%%%%%%%%%%%%%%%%%%%%%%%%%%%%%
% packages and shortcuts
\input{myPackages_gr}
\usepackage[hidelinks, linkcolor=red]{hyperref}
\input{myShortcuts}
\setlength{\parindent}{0pt}
\setlength{\parskip}{6pt}

%% exam class customization
\input{exam_config}
% ECE 444 document fonts.
%
% Body text is Barlow (SIL OFL; the faces and the license are in
% latex/fonts/Barlow/). Barlow is not in TeX Live, so it is loaded as an
% OpenType face through fontspec -- which means the document must be built
% with lualatex, not pdflatex. build_practice.sh and build_project.sh do; the
% lab packets build in ECE444-faculty from its vendored copy of this file.
%
% Compiling with pdflatex anyway is not an error: the guard below falls back
% to the LaTeX default so the source still compiles under any engine and can
% be syntax-checked. The resulting PDF just is not the one we ship.
%
% Paths are relative to latex/, which is where the build scripts run
% (each does cd "$(dirname "$0")"). Loading by path rather than by installed
% family name means no system font install and no fontconfig, so a fresh
% clone builds a byte-comparable PDF anywhere.
%
% MATH IS DELIBERATELY LEFT IN COMPUTER MODERN. Barlow has no math companion,
% and CM math under sans body text is already the house pairing: the decks set
% words in Source Sans Pro and let MathJax render symbols in its CM-derived
% face. Moving math to a sans as well would mean unicode-math plus a sans math
% font (Fira Math being the plausible one) and re-checking every equation in
% the course for symbol coverage. That is a separate decision, not a side
% effect of this one.

\usepackage{iftex}
\ifLuaTeX
	\usepackage{fontspec}
	\defaultfontfeatures{Ligatures=TeX}
	\setmainfont{Barlow}[
		Path           = fonts/Barlow/,
		Extension      = .ttf,
		UprightFont    = *-Regular,
		ItalicFont     = *-Italic,
		BoldFont       = *-Bold,
		BoldItalicFont = *-BoldItalic,
	]
	\setsansfont{Barlow}[
		Path           = fonts/Barlow/,
		Extension      = .ttf,
		UprightFont    = *-Regular,
		ItalicFont     = *-Italic,
		BoldFont       = *-SemiBold,
		BoldItalicFont = *-BoldItalic,
	]
	% microtype earns its place here specifically because of Barlow. Its
	% glyphs are narrower than Computer Modern's, so justified prose sets
	% tighter and TeX has to squeeze interword space to fit; two paragraphs
	% in the L02/L03 sets went overfull on the first Barlow build. Protrusion
	% and font expansion absorb that -- both are lualatex features, so this
	% belongs inside the guard with everything else.
	\usepackage{microtype}

	% Mono stays Latin Modern: Barlow ships no monospaced face, and verbatim
	% and code need one that actually is fixed-pitch.
\fi
   % Barlow body text (needs lualatex)
\printselectedfalse      % full worked solutions in the key, not answers-only
\CorrectChoiceEmphasis{\color{red}}
\SolutionEmphasis{\color{red}}
\renewcommand\questionlabel{\textbf{\thequestion.}}
\renewcommand{\thepartno}{\alph{partno}}
\renewcommand{\partlabel}{(\thepartno)}

% Fillable table cell: blank in the student copy, the expected value in red on
% the KEY. Fixed-width centered columns keep the blank cells writable.
\newcolumntype{K}[1]{>{\centering\arraybackslash}p{#1}}
\newcommand{\kv}[1]{\ifkey{\color{red}#1}\fi\rule[-0.10in]{0pt}{0.34in}}

% print the answers if this is the key
\ifkey
	\printanswers % if this IS the key
\else
	\noprintanswers  %if this is NOT the key
\fi

% print the header and footer
% need to print key on top if this is the key
\ifkey
	\firstpageheader{}{\color{red}{ KEY}}{}
	\runningheader{\thisclass\hspace{1pt} -- \thisterm }{L\labLesson\ Lab\hspace{1pt}\color{red}{ KEY}}{\thepage\hspace{1pt} of\hspace{1pt} \numpages}
	\firstpagefooter{}{}{}
	\runningfooter{}{}{}
\else
	\firstpageheader{}{}{}
	\runningheader{\thisclass \hspace{1pt} -- \thisterm }{L\labLesson\ Lab\hspace{1pt}}{\thepage\hspace{1pt} of\hspace{1pt} \numpages}
	\firstpagefooter{}{}{}
	\runningfooter{}{}{}
\fi

\runningheadrule


\begin{document}
\pagestyle{headandfoot}   % exam class


%%%%%%%%%%%%%%%%%%%%%%%%%%%%%%%%%%%%%%%%%%%%%%%%%%%%%%
% title block

\begin{center}
USAF ACADEMY \\
DEPARTMENT OF ELECTRICAL AND COMPUTER ENGINEERING \\
LESSON \labLesson\ LAB - Antenna Pattern \\
\textbf{Lab Packet} \\
\thisterm \\[4mm]
\end{center}

\vspace{-4pt}
\textbf{Name:} \rule[-2pt]{2.6in}{0.4pt} \hfill
\textbf{Section:} \rule[-2pt]{0.9in}{0.4pt} \hfill
\textbf{Date:} \rule[-2pt]{1.1in}{0.4pt}

\vspace{4pt}

\fullwidth{\textbf{Learning Objective 3.6: } I can distinguish between array factor
and true antenna pattern and account for element pattern effects.}

\vspace{2pt}

This packet is the L\labLesson\ Lab, in three parts, and you turn it in once,
complete, before Lesson 22. Do \textbf{Part 1} on your own before the lab
period: the predictions you will test, and the error budgets of the two ways you
will measure the pattern. You need it at the bench. \textbf{Part 2} is your own
measurements, from your own kit, recorded in class. Do \textbf{Part 3} on your
own after the lab: reconcile the measurements with the predictions, compare the
two methods, and work two analysis problems.

\textbf{Grading.} One submission, counted toward your engagement grade and
graded as a genuine, documented attempt.

The array is the PHASER's: $N = 8$ elements, E1 to E8, on a $d = 14\ \text{mm}$
pitch, so the aperture spans $D = 98\ \text{mm}$ between the outer element
centers. Facing the patches, E1 is on the right and E8 on the left, and
$+\theta$ is on the E8 side. Use $c = 3.00\times 10^{8}\ \m/\text{s}$, and report
every lobe level relative to its own trace's main-lobe peak, in dBc, unless a
question asks otherwise.

\begin{questions}

\fullwidth{\large\textbf{Part 1: Pre-Lab} \normalsize(individual, before the lab period)}

%%%%%%%%%%%%%%%%%%%%%%%%%%%%%%%%%%%%%%%%%%%%%%%%%%%
\question \textbf{Frequency, range, and distance.} Use the HB100 frequency your
kit reported with \textbf{Find HB100} in the L18 Lab. If you did not record one,
use the nominal $10.525\ \ghz$ and say so.

\begin{parts}

	\begin{minipage}[t]{0.58\linewidth}
		\part Record the frequency you are using, then find the wavelength, the
		electrical spacing $d/\lambda$, and $kd$ in degrees.
		\begin{solution}[1.5in]
		The key uses the nominal $10.525\ \ghz$:
		\eq{ \lambda &= \frac{c}{f} = \frac{3.00\times 10^{8}}{10.525\times 10^{9}} = 28.50\ \text{mm} \\
		     \frac{d}{\lambda} &= \frac{14}{28.50} = 0.491 \\
		     kd &= 360\mydeg\ \frac{d}{\lambda} = 176.8\mydeg }
		\end{solution}
	\end{minipage}\hfill%
	\begin{minipage}[t]{0.37\linewidth}
		\raggedleft \vspace{12pt}
		$f$ = \ansbox[1.3in]{$10.525\ \ghz$}\\[10pt]
		$d/\lambda$ = \ansbox[1.3in]{$0.491$}\\[10pt]
		$kd$ = \ansbox[1.3in]{$176.8\mydeg$}
	\end{minipage}

	\begin{minipage}[t]{0.58\linewidth}
		\part Find the far-field distance of the array at your frequency.
		\begin{solution}[1.2in]
		\eq{ r \ge \frac{2D^2}{\lambda} = \frac{2\lp0.098\ \m\rp^2}{0.02850\ \m} = 0.67\ \m }
		\end{solution}
	\end{minipage}\hfill%
	\begin{minipage}[t]{0.37\linewidth}
		\raggedleft \vspace{12pt}
		$r$ = \ansbox[1.3in]{$0.67\ \m$}
	\end{minipage}

	\part Today's source sits $1.2\ \m$ from the array, nearly twice the
	far-field distance. In the Module 3 dry run, the drop in peak level at
	$30\mydeg$ read $-0.7\ \db$ with the source at $1.2\ \m$ and only $-0.2$ to
	$-0.4\ \db$ at $1\ \m$, against an ideal-element prediction of
	$-0.6\ \db$. In two or three sentences, explain why a distance that clears
	the far-field criterion can still be too close for a level comparison of
	under a decibel.
		\begin{solution}[1.4in]
		The criterion $2D^2/\lambda$ limits the phase curvature across the
		aperture to $22.5\mydeg$, which protects the pattern's shape: its angles,
		widths, and nulls. A level comparison of a few tenths of a decibel is a
		much finer test, and at $1\ \m$ the residual curvature, the reflections
		from the bench, and the near-field coupling of the source's own mount all
		contribute errors of that size. Moving the source out reduces each of
		them, which is why the dry run read the expected drop only at
		$1.2\ \m$.
		\end{solution}

\end{parts}

%%%%%%%%%%%%%%%%%%%%%%%%%%%%%%%%%%%%%%%%%%%%%%%%%%%
\question \textbf{Reading dBc off a trace.} A beam-held trace gives peak
readings of $-18.4\ \text{dBFS}$ for the main lobe, $-31.2\ \text{dBFS}$ for the
left first sidelobe, and $-29.9\ \text{dBFS}$ for the right first sidelobe.

\begin{parts}

	\begin{minipage}[t]{0.58\linewidth}
		\part Convert both sidelobe readings to dBc.
		\begin{solution}[1.1in]
		A dBc value is the lobe peak referred to the main-lobe peak:
		\eq{ L_\text{left} &= -31.2 - (-18.4) = -12.8\ \text{dBc} \\
		     L_\text{right} &= -29.9 - (-18.4) = -11.5\ \text{dBc} }
		\end{solution}
	\end{minipage}\hfill%
	\begin{minipage}[t]{0.37\linewidth}
		\raggedleft \vspace{12pt}
		$L_\text{left}$ = \ansbox[1.25in]{$-12.8$ dBc}\\[12pt]
		$L_\text{right}$ = \ansbox[1.25in]{$-11.5$ dBc}
	\end{minipage}

	\part Raising \textbf{Rx Gain} by $6\ \db$ and repeating the measurement
	changes the three dBFS readings. State what happens to them and to the two
	dBc values, and explain why the two answers differ.
		\begin{solution}[1.3in]
		All three dBFS readings rise by about $6\ \db$, because the gain change
		scales every sample by the same factor. The dBc values do not change,
		because each is the difference of two readings taken through the same
		gain, and the common term cancels. That is why a pattern is reported
		relative to its own peak: the result does not depend on the receiver's
		absolute scale, which is not calibrated here. It holds only while the
		receiver stays linear, which is why Part 2 sets Rx Gain for a peak near
		$-10\ \text{dBFS}$ rather than near full scale.
		\end{solution}

\end{parts}

%%%%%%%%%%%%%%%%%%%%%%%%%%%%%%%%%%%%%%%%%%%%%%%%%%%
\newpage
\question\label{q:broad} \textbf{The beam held at broadside.} In Part 2 you will hold the beam
at $0\mydeg$ and measure the pattern with the source at known angles. Predict
what you should see. The array factor of a uniform array is
$\vert AF\vert = \vert\sin(N\psi/2)/(N\sin(\psi/2))\vert$ with
$\psi = kd\lp\sin\theta - \sin\theta_0\rp$, and the ideal element's power
pattern is $\cos\theta$ (Lesson 21).

\begin{parts}

	\begin{minipage}[t]{0.58\linewidth}
		\part Find the angles of the first three nulls on each side.
		\begin{solution}[1.5in]
		The nulls fall where $N\psi/2 = m\pi$, so
		\eq{ \sin\theta_m &= \pm\frac{m\lambda}{Nd} = \pm m\lp0.2545\rp \\
		     \theta_1, \theta_2, \theta_3 &= \pm 14.7\mydeg,\ \pm 30.6\mydeg,\ \pm 49.8\mydeg }
		\end{solution}
	\end{minipage}\hfill%
	\begin{minipage}[t]{0.37\linewidth}
		\raggedleft \vspace{12pt}
		$\theta_1$ = \ansbox[1.3in]{$\pm 14.7\mydeg$}\\[10pt]
		$\theta_2$ = \ansbox[1.3in]{$\pm 30.6\mydeg$}\\[10pt]
		$\theta_3$ = \ansbox[1.3in]{$\pm 49.8\mydeg$}
	\end{minipage}

	\begin{minipage}[t]{0.58\linewidth}
		\part The first sidelobes of the array factor sit at $\pm 21.4\mydeg$,
		$12.8\ \db$ below the peak. Include the element pattern and give the
		first-sidelobe level you expect to measure, in dBc.
		\begin{solution}[1.3in]
		The element pattern lowers each sidelobe by its own roll-off, and the main
		lobe at $0\mydeg$ not at all:
		\eq{ 10\mylog{\cosp{21.4\mydeg}} &= -0.3\ \db \\
		     L_1 &= -12.8 - 0.3 = -13.1\ \text{dBc} }
		\end{solution}
	\end{minipage}\hfill%
	\begin{minipage}[t]{0.37\linewidth}
		\raggedleft \vspace{12pt}
		$L_1$ = \ansbox[1.3in]{$-13.1$ dBc}
	\end{minipage}

	\part Part 2 records the level at $10\mydeg$ marks across the front of the
	array. Using your null angles, name the marks that fall within a few
	degrees of a null, and explain in one or two sentences why the coarse marks
	alone cannot give you the sidelobe levels.
		\begin{solution}[1.3in]
		The $\pm 30\mydeg$ marks sit $0.6\mydeg$ from the second nulls and the
		$\pm 50\mydeg$ marks $0.2\mydeg$ from the third, so those four read near
		the floor. The first sidelobes peak at $\pm 21.4\mydeg$, between the
		$20\mydeg$ and $30\mydeg$ marks, and a lobe about $15\mydeg$ wide changes
		by several decibels over a few degrees, so a fixed mark rarely lands on its
		peak. That is why Part 2 has you turn the kit slowly through each lobe and
		record its highest reading.
		\end{solution}

\end{parts}

%%%%%%%%%%%%%%%%%%%%%%%%%%%%%%%%%%%%%%%%%%%%%%%%%%%
\newpage
\question\label{q:steer} \textbf{The beam held at $30\mydeg$.} Now hold the beam at
$\theta_0 = 30\mydeg$.

\begin{parts}

	\begin{minipage}[t]{0.58\linewidth}
		\part Find the first null on each side of the steered beam.
		\begin{solution}[1.3in]
		\eq{ \sin\theta &= \sin 30\mydeg \pm 0.2545 = 0.2455,\ 0.7545 \\
		     \theta &= 14.2\mydeg,\ 49.0\mydeg }
		The pattern is no longer symmetric: the inner null sits $15.8\mydeg$
		below the commanded angle and the outer null $19.0\mydeg$ above it.
		\end{solution}
	\end{minipage}\hfill%
	\begin{minipage}[t]{0.37\linewidth}
		\raggedleft \vspace{12pt}
		inner = \ansbox[1.3in]{$14.2\mydeg$}\\[10pt]
		outer = \ansbox[1.3in]{$49.0\mydeg$}
	\end{minipage}

	\begin{minipage}[t]{0.58\linewidth}
		\part Find the drop in peak level relative to the broadside beam of
		Question \ref{q:broad}, for the ideal element and for a patch that rolls off as
		$\cos^{1.4}\theta$ in power.
		\begin{solution}[1.3in]
		\eq{ \Delta G_\text{ideal} &= 10\mylog{\cosp{30\mydeg}} = -0.6\ \db \\
		     \Delta G_\text{patch} &= 14\mylog{\cosp{30\mydeg}} = -0.9\ \db }
		\end{solution}
	\end{minipage}\hfill%
	\begin{minipage}[t]{0.37\linewidth}
		\raggedleft \vspace{12pt}
		ideal = \ansbox[1.3in]{$-0.6\ \db$}\\[10pt]
		patch = \ansbox[1.3in]{$-0.9\ \db$}
	\end{minipage}

	\part The array factor alone puts both first sidelobes $12.8\ \db$ below the
	peak, the inner one at $7.7\mydeg$ and the outer one at $60.0\mydeg$. Fill
	in the table for the array factor alone and for the element pattern times
	the array factor. Express each sidelobe relative to its own trace's peak.

	\begin{center}
	\renewcommand{\arraystretch}{1.25}
	\begin{tabular}{| l | K{1.4in} | K{1.4in} |}
	\hline
	& \textbf{Array factor alone} & \textbf{Element $\times$ array} \\\hline\hline
	Peak angle & \kv{$30.0\mydeg$} & \kv{$29.6\mydeg$} \\\hline
	Inner first sidelobe & \kv{$-12.8$ dBc} & \kv{$-12.2$ dBc} \\\hline
	Outer first sidelobe & \kv{$-12.8$ dBc} & \kv{$-15.1$ dBc} \\\hline
	\end{tabular}
	\end{center}

		\begin{solution}[1.6in]
		The element pattern lowers the peak by $0.6\ \db$ and each sidelobe by
		its own roll-off. The inner lobe sits where the element is nearly flat,
		$10\mylog{\cosp{7.7\mydeg}} = -0.04\ \db$, so referred to the lowered peak
		it reads $-12.8 + 0.6 = -12.2$ dBc. The outer lobe loses
		$10\mylog{\cosp{58.8\mydeg}} = -2.9\ \db$, so it reads
		$-12.8 + 0.6 - 2.9 = -15.1$ dBc. The element pattern also pulls the peak
		$0.4\mydeg$ toward broadside and the outer lobe from $60.0\mydeg$ to
		$58.8\mydeg$.
		\end{solution}

	\part Part 2 also takes one reading with the kit turned to $180\mydeg$, the
	source behind the board. Predict that reading, and give the reason.
		\begin{solution}[1.0in]
		It should read at the floor. The ground plane under the patches removes
		the back hemisphere, so the element pattern, and with it the antenna
		pattern, is zero behind the board whatever the array factor does there.
		\end{solution}

\end{parts}

%%%%%%%%%%%%%%%%%%%%%%%%%%%%%%%%%%%%%%%%%%%%%%%%%%%
\newpage
\question\label{q:budget} \textbf{Two ways to move the source, and their error budgets.} Part 2
measures the pattern twice with the beam held fixed. In the \textbf{hand walk}
you carry the HB100 around a taped arc of radius $1.2\ \m$ while the GUI plots
received level against time. In the \textbf{rotation} the HB100 stays on its
tripod $1.2\ \m$ out, and you turn the kit about the center of the array over a
protractor, reading the level at known angles.

\begin{parts}

	\begin{minipage}[t]{0.58\linewidth}
		\part During a walk your radius drifts by $10\ \text{cm}$ either way.
		Received power goes as $1/r^2$. Find the change in the reading for a
		radius of $1.1\ \m$ and of $1.3\ \m$.
		\begin{solution}[1.3in]
		\eq{ 20\mylog{\frac{1.2}{1.1}} &= +0.8\ \db \\
		     20\mylog{\frac{1.2}{1.3}} &= -0.7\ \db }
		Either drift is as large as the $0.6\ \db$ scan loss you want to measure.
		\end{solution}
	\end{minipage}\hfill%
	\begin{minipage}[t]{0.37\linewidth}
		\raggedleft \vspace{12pt}
		$1.1\ \m$: \ansbox[1.2in]{$+0.8\ \db$}\\[10pt]
		$1.3\ \m$: \ansbox[1.2in]{$-0.7\ \db$}
	\end{minipage}

	\begin{minipage}[t]{0.58\linewidth}
		\part The walk's x-axis is time. If you convert it to angle by assuming
		a constant walking speed, and your hand actually moves up to $30\%$
		faster or slower than average while crossing the main lobe, give the
		range of half-power beamwidths the conversion reports for a true
		$13.0\mydeg$ beam.
		\begin{solution}[1.3in]
		A fast hand spends less time in the lobe, so the constant-speed map draws
		it narrower; a slow hand draws it wider:
		\eq{ \theta_\text{fast} &= \frac{13.0\mydeg}{1.30} = 10.0\mydeg \\
		     \theta_\text{slow} &= \frac{13.0\mydeg}{0.70} = 18.6\mydeg }
		\end{solution}
	\end{minipage}\hfill%
	\begin{minipage}[t]{0.37\linewidth}
		\raggedleft \vspace{12pt}
		range = \ansbox[1.4in]{$10.0\mydeg$ to $18.6\mydeg$}
	\end{minipage}

	\begin{minipage}[t]{0.58\linewidth}
		\part In the rotation, suppose the kit turns about a point $20\ \text{mm}$
		from the center of the array. Find the largest angle error and the
		largest change in level this offset can produce at $1.2\ \m$.
		\begin{solution}[1.4in]
		As the kit turns, the array center moves on a $20\ \text{mm}$ circle.
		Moved sideways, it shifts the source's angle by up to
		\eq{ \Delta\theta = \arctan\frac{0.020}{1.2} = 1.0\mydeg }
		and moved along the line to the source, it changes the distance by up to
		$20\ \text{mm}$:
		\eq{ 20\mylog{\frac{1.22}{1.20}} = 0.14\ \db }
		\end{solution}
	\end{minipage}\hfill%
	\begin{minipage}[t]{0.37\linewidth}
		\raggedleft \vspace{12pt}
		$\Delta\theta$ = \ansbox[1.2in]{$1.0\mydeg$}\\[10pt]
		$\Delta L$ = \ansbox[1.2in]{$0.14\ \db$}
	\end{minipage}

	\part A reflection from the room adds to the direct path in both methods.
	Find the relative amplitude $\rho$ of a single reflected path that moves the
	reading by $+1\ \db$ when the two add in phase, and express $\rho$ in dB.
		\begin{solution}[1.2in]
		\eq{ 20\mylog{1+\rho} = 1.0\ \db \ \rightarrow\ \rho = 10^{0.05} - 1 = 0.122 = -18.3\ \db }
		A reflection only $18\ \db$ below the direct path moves a reading by a
		decibel, so the room contributes errors of the same size as the effects
		this lab measures.
		\end{solution}

	\part In two or three sentences, say which of the errors above the rotation
	removes, and which it leaves in place.
		\begin{solution}[1.3in]
		The rotation removes the radius drift, because the source never moves,
		and the time-for-angle error, because each reading is taken at a marked
		angle; it also keeps the HB100 aimed at the array without effort. It
		leaves the reflections, which depend on the room rather than on how the
		source moves, and it adds the axis offset and the protractor's reading
		error, both of which are small and can be measured.
		\end{solution}

\end{parts}

%%%%%%%%%%%%%%%%%%%%%%%%%%%%%%%%%%%%%%%%%%%%%%%%%%%
\vspace{12pt}
\fullwidth{\large\textbf{Part 2: Measurements} \normalsize(your kit, in class)}

Follow the procedure on the L\labLesson\ Lab page. Before you record anything,
set \textbf{Rx Gain} so that the main-lobe peak, with the source at boresight
and the beam at $0\mydeg$, reads $-10 \pm 3\ \text{dBFS}$ on the FFT, and leave
it there for the whole lab. Keep the HB100 at patch-row height, $1.2\ \m$ from
the array. Record each reading as you take it.

Signal Freq from \textbf{Find HB100}: \rule[-2pt]{1.3in}{0.4pt} \hfill
Rx Gain: \rule[-2pt]{0.8in}{0.4pt} \hfill
Source distance: \rule[-2pt]{0.8in}{0.4pt}

%%%%%%%%%%%%%%%%%%%%%%%%%%%%%%%%%%%%%%%%%%%%%%%%%%%
\question \textbf{Table A: check the angle axis.} With the beam sweeping as in
the L18 Lab, turn the kit so the source sits at each protractor angle below,
and record \textbf{Est. Angle}. Turning the kit's face toward E1, as seen from
the source, puts the source on the $+$ side.

\begin{center}
\renewcommand{\arraystretch}{1.25}
\begin{tabular}{| l | K{1.0in} | K{1.0in} | K{1.0in} |}
\hline
Protractor angle & $-30\mydeg$ & $0\mydeg$ & $+30\mydeg$ \\\hline
Est. Angle & \kv{$-30 \pm 2\mydeg$} & \kv{$0 \pm 2\mydeg$} & \kv{$+30 \pm 2\mydeg$} \\\hline
\end{tabular}
\end{center}

%%%%%%%%%%%%%%%%%%%%%%%%%%%%%%%%%%%%%%%%%%%%%%%%%%%
\question \textbf{Table B: the electrical sweep.} Kit back at $0\mydeg$, beam
sweeping. Read the trace.

\begin{center}
\renewcommand{\arraystretch}{1.25}
\begin{tabular}{| K{1.1in} | K{1.0in} | K{1.2in} | K{1.2in} |}
\hline
\textbf{Peak (dBFS)} & \textbf{HPBW} & \textbf{Left sidelobe (dBc)} & \textbf{Right sidelobe (dBc)} \\\hline
\kv{$-10 \pm 3$} & \kv{$13.0 \pm 1\mydeg$} & \kv{$-11$ to $-13$} & \kv{$-11$ to $-13$} \\\hline
\end{tabular}
\end{center}

%%%%%%%%%%%%%%%%%%%%%%%%%%%%%%%%%%%%%%%%%%%%%%%%%%%
\question \textbf{Table C: the hand walk.} Beam held at $0\mydeg$. Walk the
HB100 from $-90\mydeg$ to $+90\mydeg$ on the taped $1.2\ \m$ arc, aimed at the
array, three times. Record each walk's main-lobe peak and both first sidelobes.

\begin{center}
\renewcommand{\arraystretch}{1.25}
\begin{tabular}{| l | K{1.0in} | K{1.0in} | K{1.0in} | K{1.0in} |}
\hline
& \textbf{Walk time (s)} & \textbf{Peak (dBFS)} & \textbf{Left SL (dBc)} & \textbf{Right SL (dBc)} \\\hline\hline
Walk 1 & \kv{$7$ to $15$} & \kv{$-10 \pm 3$} & \kv{$-10$ to $-16$} & \kv{$-10$ to $-16$} \\\hline
Walk 2 & & & & \\\hline
Walk 3 & & & & \\\hline
Spread (max $-$ min) & & \kv{$\sim 1\ \db$} & \kv{$1$ to $3\ \db$} & \kv{$1$ to $3\ \db$} \\\hline
\end{tabular}
\end{center}

%%%%%%%%%%%%%%%%%%%%%%%%%%%%%%%%%%%%%%%%%%%%%%%%%%%
\newpage
\question \textbf{Table D: the rotation, beam held at $0\mydeg$.} HB100 on its
tripod at boresight, $1.2\ \m$ out. Turn the kit to each mark and record the
level relative to the reading at $0\mydeg$, in dB. Write ``floor'' where the
reading stops changing as you turn.

\begin{center}
\small
\setlength{\tabcolsep}{3pt}
\renewcommand{\arraystretch}{1.25}
\begin{tabular}{| l | K{0.48in} | K{0.48in} | K{0.48in} | K{0.48in} | K{0.48in} | K{0.48in} | K{0.48in} | K{0.48in} | K{0.48in} | K{0.48in} |}
\hline
Angle & $-90\mydeg$ & $-80\mydeg$ & $-70\mydeg$ & $-60\mydeg$ & $-50\mydeg$ & $-40\mydeg$ & $-30\mydeg$ & $-20\mydeg$ & $-10\mydeg$ & $0\mydeg$ \\\hline
Level & \kv{floor} & \kv{floor} & \kv{$-24$} & \kv{$-21$} & \kv{floor} & \kv{$-18$} & \kv{floor} & \kv{$-14$} & \kv{$-8$} & \kv{$0$} \\\hline\hline
Angle & $+10\mydeg$ & $+20\mydeg$ & $+30\mydeg$ & $+40\mydeg$ & $+50\mydeg$ & $+60\mydeg$ & $+70\mydeg$ & $+80\mydeg$ & $+90\mydeg$ & \\\hline
Level & \kv{$-8$} & \kv{$-14$} & \kv{floor} & \kv{$-18$} & \kv{floor} & \kv{$-21$} & \kv{$-24$} & \kv{floor} & \kv{floor} & \\\hline
\end{tabular}
\end{center}

Now turn the kit slowly through each lobe and null below, and record the angle
and level of its extreme, the highest reading for a lobe and the lowest for a
null.

\begin{center}
\renewcommand{\arraystretch}{1.25}
\begin{tabular}{| l | K{1.2in} | K{1.3in} |}
\hline
& \textbf{Angle} & \textbf{Level} \\\hline\hline
Main lobe & \kv{$0 \pm 1\mydeg$} & \kv{$-10 \pm 3\ \text{dBFS}$} \\\hline
Left first null & \kv{$-14.7 \pm 1.5\mydeg$} & \kv{floor} \\\hline
Right first null & \kv{$+14.7 \pm 1.5\mydeg$} & \kv{floor} \\\hline
Left first sidelobe & \kv{$-21.4 \pm 2\mydeg$} & \kv{$-13.1 \pm 2$ dBc} \\\hline
Right first sidelobe & \kv{$+21.4 \pm 2\mydeg$} & \kv{$-13.1 \pm 2$ dBc} \\\hline
\end{tabular}
\end{center}

%%%%%%%%%%%%%%%%%%%%%%%%%%%%%%%%%%%%%%%%%%%%%%%%%%%
\newpage
\question \textbf{Table E: the rotation, beam held at $30\mydeg$.} Same setup,
beam held at $30\mydeg$. Record each level relative to the \emph{Table D}
main-lobe peak, so that the drop in peak level shows directly.

\begin{center}
\small
\setlength{\tabcolsep}{3pt}
\renewcommand{\arraystretch}{1.25}
\begin{tabular}{| l | K{0.48in} | K{0.48in} | K{0.48in} | K{0.48in} | K{0.48in} | K{0.48in} | K{0.48in} | K{0.48in} | K{0.48in} | K{0.48in} |}
\hline
Angle & $-90\mydeg$ & $-80\mydeg$ & $-70\mydeg$ & $-60\mydeg$ & $-50\mydeg$ & $-40\mydeg$ & $-30\mydeg$ & $-20\mydeg$ & $-10\mydeg$ & $0\mydeg$ \\\hline
Level & \kv{floor} & \kv{floor} & \kv{$-22$} & \kv{$-21$} & \kv{floor} & \kv{$-19$} & \kv{floor} & \kv{$-20$} & \kv{$-18$} & \kv{floor} \\\hline\hline
Angle & $+10\mydeg$ & $+20\mydeg$ & $+30\mydeg$ & $+40\mydeg$ & $+50\mydeg$ & $+60\mydeg$ & $+70\mydeg$ & $+80\mydeg$ & $+90\mydeg$ & $180\mydeg$ \\\hline
Level & \kv{$-14$} & \kv{$-7$} & \kv{$-0.6$} & \kv{$-6$} & \kv{floor} & \kv{$-16$} & \kv{$-21$} & \kv{floor} & \kv{floor} & \kv{floor} \\\hline
\end{tabular}
\end{center}

\begin{center}
\renewcommand{\arraystretch}{1.25}
\begin{tabular}{| l | K{1.2in} | K{1.3in} |}
\hline
& \textbf{Angle} & \textbf{Level} \\\hline\hline
Main lobe & \kv{$29.6 \pm 1\mydeg$} & \kv{$-0.6$ to $-0.9\ \db$ rel.\ Table D} \\\hline
Inner first null & \kv{$14.2 \pm 1.5\mydeg$} & \kv{floor} \\\hline
Outer first null & \kv{$49.0 \pm 1.5\mydeg$} & \kv{floor} \\\hline
Inner first sidelobe & \kv{$7.7 \pm 2\mydeg$} & \kv{$-12.2 \pm 2$ dBc} \\\hline
Outer first sidelobe & \kv{$58.8 \pm 2\mydeg$} & \kv{$-15.1 \pm 2$ dBc} \\\hline
\end{tabular}
\end{center}

\begin{solution}
The coarse marks are the element-times-array prediction at $10.525\ \ghz$,
rounded to the decibel; ``floor'' marks the entries more than about $25\ \db$
down, where the room and the receiver set the reading. Sidelobes in the lobe
tables are relative to the trace's own peak. Expect the measured floor near
$-23$ to $-30$ dB, and expect individual readings to scatter by $\pm 1\ \db$
from reflections even though the rotation repeats well.
\end{solution}

%%%%%%%%%%%%%%%%%%%%%%%%%%%%%%%%%%%%%%%%%%%%%%%%%%%
\newpage
\fullwidth{\large\textbf{Part 3: Writeup} \normalsize(individual, after the lab)}

%%%%%%%%%%%%%%%%%%%%%%%%%%%%%%%%%%%%%%%%%%%%%%%%%%%
\question \textbf{Reconcile the rotation with the two predictions.} Answer each
part in three or four sentences, using your numbers from Part 2 and your own
predictions from Part 1.

\begin{parts}

	\part Compare the first sidelobes of Table D with your prediction from
	Question \ref{q:broad}, and with the electrical sweep of Table B. Say whether the
	element pattern is visible at broadside.
		\begin{solutionorlines}[1.6in]
		Both sidelobes should land within about $2\ \db$ of $-13.1$ dBc, and the
		sweep's within the same window of $-12.8$ dBc. At $\pm 21\mydeg$ the
		element pattern changes the level by only $0.3\ \db$, well inside the
		room's $\pm 1\ \db$, so at broadside the element pattern is not
		distinguishable and the array factor alone describes the pattern. Two
		sidelobes that differ by a decibel or so are normal in a room; a
		difference of several decibels points to a calibration problem.
		\end{solutionorlines}

	\part Find the drop in peak level from Table D to Table E and compare it
	with the ideal and patch predictions of Question \ref{q:steer}. State whether your
	measurement resolves the drop, given the error budget of Question \ref{q:budget}.
		\begin{solutionorlines}[1.6in]
		The dry run measured $-0.7\ \db$, between the ideal $-0.6\ \db$ and the
		patch's $-0.9\ \db$. Because the source never moved, the radius term of
		the budget is gone and the axis offset contributes at most
		$0.14\ \db$; the two peaks are also taken at nearly the same physical
		geometry, so a reflection biases both by a similar amount. A drop of
		$0.5$ to $1.0\ \db$ that repeats on a second rotation resolves the scan
		loss; a single reading does not, by the $\pm 1\ \db$ reflection term.
		\end{solutionorlines}

	\part Compare the inner and outer first sidelobes of Table E with the two
	columns of your Question \ref{q:steer} table. Which model does your measurement
	support, and how large is the asymmetry compared with the measurement's
	uncertainty?
		\begin{solutionorlines}[1.6in]
		The array factor alone predicts equal sidelobes; element times array
		predicts $-12.2$ and $-15.1$ dBc, a $2.9\ \db$ asymmetry with the outer
		lobe lower. A measured asymmetry of $2$ to $4\ \db$ in that direction
		supports the product, and it exceeds the $\pm 2\ \db$ worst-case
		uncertainty on a difference of two readings only when it repeats. An
		asymmetry in the other direction would point to a reflection near
		$60\mydeg$ or to the kit turned the wrong way.
		\end{solutionorlines}

	\part Report your $180\mydeg$ reading and explain it. Then explain why the
	electrical sweep of Table B cannot show the sidelobe asymmetry of Table E,
	however carefully it is taken.
		\begin{solutionorlines}[1.8in]
		The $180\mydeg$ reading sits at the floor, because the ground plane
		removes the back hemisphere. A beam sweep holds the source at one
		physical angle, so every point of the trace sees the same element gain:
		the element pattern multiplies the whole trace by one constant and drops
		out of every dBc value. Only a measurement that moves the source across
		the element pattern, the walk or the rotation, records the element
		factor, which is the difference between the array factor and the
		antenna pattern.
		\end{solutionorlines}

\end{parts}

%%%%%%%%%%%%%%%%%%%%%%%%%%%%%%%%%%%%%%%%%%%%%%%%%%%
\newpage
\question \textbf{The walk against the rotation.} Use Table C and Tables D and E.

\begin{parts}

	\part Compare the spread of your three walks with the repeatability you saw
	in the rotation. Could a single walk resolve the $0.6\ \db$ scan loss? The
	$2.9\ \db$ sidelobe asymmetry? Give the reason for each answer.
		\begin{solutionorlines}[2.0in]
		Three walks typically spread by about a decibel in peak level and one to
		three in the sidelobes, from radius drift, aim, and walking speed, while a
		repeated rotation reading agrees to a few tenths. A single walk cannot
		resolve $0.6\ \db$, because a $10\ \text{cm}$ radius drift alone moves
		the reading by $0.7$ to $0.8\ \db$. The $2.9\ \db$ asymmetry is larger
		than the walk's scatter, so a careful walk can show its direction, but
		not its size to better than a decibel or two.
		\end{solutionorlines}

	\begin{minipage}[t]{0.58\linewidth}
		\part A walk runs from $t = 0.00\ \text{s}$ at $-90\mydeg$ to
		$t = 7.40\ \text{s}$ at $+90\mydeg$, with the beam held at $0\mydeg$, and
		the main-lobe peak lands at $t = 3.10\ \text{s}$. Convert that time to an
		angle assuming a constant walking speed, and state the error.
		\begin{solution}[1.2in]
		\eq{ \theta &= -90\mydeg + 180\mydeg\lp\frac{3.10}{7.40}\rp \\
		            &= -14.6\mydeg }
		The peak is at $0\mydeg$, so the constant-speed map is $14.6\mydeg$ off.
		\end{solution}
	\end{minipage}\hfill%
	\begin{minipage}[t]{0.37\linewidth}
		\raggedleft \vspace{12pt}
		$\theta$ = \ansbox[1.25in]{$-14.6\mydeg$}\\[12pt]
		error = \ansbox[1.25in]{$14.6\mydeg$}
	\end{minipage}

	\part Explain why the lobe \emph{levels} on a walked trace are usable while
	the lobe \emph{angles} are not, and name what an antenna range uses instead
	of a protractor to record the angle.
		\begin{solutionorlines}[1.6in]
		The level at each instant is the pattern at the angle the source occupied
		then, and it does not depend on how fast the source was moving; walking
		slowly through a lobe draws it wider without changing its height. The
		angle is never recorded, and recovering it from time needs the walking
		speed at every instant. A range turns the antenna on a turntable with a
		shaft encoder, which tags every sample with its angle.
		\end{solutionorlines}

	\begin{minipage}[t]{0.58\linewidth}
		\part The main-lobe peak and a sidelobe peak each carry their own
		$\pm 1\ \db$ of reflection ripple. Give the worst-case uncertainty on the
		sidelobe level in dBc, and the range a true $-13$ dBc sidelobe could read.
		\begin{solution}[1.0in]
		A dBc value is a difference of two readings, and the two errors can have
		opposite signs: $\pm 2\ \db$, so anywhere from $-11$ to $-15$ dBc.
		\end{solution}
	\end{minipage}\hfill%
	\begin{minipage}[t]{0.37\linewidth}
		\raggedleft \vspace{12pt}
		uncertainty = \ansbox[1.1in]{$\pm 2\ \db$}\\[12pt]
		range = \ansbox[1.25in]{$-11$ to $-15$ dBc}
	\end{minipage}

	\part A customer wants written proof that this array's sidelobes are below
	$-25$ dBc. State whether either of today's measurements can support that
	claim, and give the two reasons that decide it.
		\begin{solutionorlines}[1.4in]
		Neither can. A $\pm 2\ \db$ window on every dBc value cannot certify a level
		below a fixed number, since a measured $-26$ dBc could be a true $-24$. And
		the floor sits about $23$ to $30\ \db$ below the peak, so a $-25$ dBc
		sidelobe may not rise out of it at all. The claim needs an anechoic
		chamber, which removes the reflections, and a receiver with more dynamic
		range.
		\end{solutionorlines}

\end{parts}

%%%%%%%%%%%%%%%%%%%%%%%%%%%%%%%%%%%%%%%%%%%%%%%%%%%
\newpage
\question \textbf{How far out is far enough.} A Ku-band panel with a
$0.200\ \m$ aperture is to be measured at $15.0\ \ghz$.

\begin{parts}

	\begin{minipage}[t]{0.58\linewidth}
		\part Find its far-field distance.
		\begin{solution}[1.2in]
		\eq{ \lambda &= \frac{3.00\times 10^8}{15.0\times 10^9} = 0.0200\ \m \\
		     r &\ge \frac{2\lp0.200\rp^2}{0.0200} = 4.00\ \m }
		Doubling the aperture and halving the wavelength multiply the distance
		by eight.
		\end{solution}
	\end{minipage}\hfill%
	\begin{minipage}[t]{0.37\linewidth}
		\raggedleft \vspace{12pt}
		$r$ = \ansbox[1.25in]{$4.00\ \m$}
	\end{minipage}

	\begin{minipage}[t]{0.58\linewidth}
		\part The panel is measured on today's $1.2\ \m$ range anyway. The extra
		path from the source to the aperture edge, relative to its center, is
		about $D^2/8r$. Find the edge phase error and compare it with the
		$22.5\mydeg$ the far-field criterion allows.
		\begin{solution}[1.3in]
		\eq{ \Delta &= \frac{\lp0.200\rp^2}{8\lp1.2\rp} = 4.2\ \text{mm} \\
		     \Delta\phi &= 360\mydeg\ \frac{4.2}{20.0} = 75\mydeg }
		That is more than three times the $22.5\mydeg$ the criterion allows.
		\end{solution}
	\end{minipage}\hfill%
	\begin{minipage}[t]{0.37\linewidth}
		\raggedleft \vspace{12pt}
		$\Delta\phi$ = \ansbox[1.25in]{$75\mydeg$}
	\end{minipage}

	\part Describe what that phase error does to the measured pattern, and say
	which pattern feature it corrupts worst.
		\begin{solutionorlines}[1.4in]
		The edges of the aperture receive the wave out of phase with its center,
		so the measurement no longer sees a plane wave and the aperture is
		defocused. The main lobe broadens and its peak drops, but the damage is
		worst in the nulls and sidelobes: a deep null depends on precise phase
		alignment across the aperture, so the nulls fill and the sidelobes read
		high. A sidelobe level measured this way is pessimistic, and a null depth
		is meaningless.
		\end{solutionorlines}

\end{parts}

\end{questions}

\vspace{1.0em}
\noindent\textbf{Documentation:} \rule[-2pt]{0.6\linewidth}{0.4pt}

\end{document}
